\documentclass[12px]{article}

\title{Lezione 13 Metodi Numerici di approssimazione}
\date{2026-04-28}
\author{Federico De Sisti}

\input{../../setup.tex}

\begin{document}
	\maketitle
	\newpage
	\subsection{RK2}
	$K_1 = f(t_n,u_n),  K_2 = f(t_n+C_2h, u_n + d_2;hk_1$
	\[
	\begin{cases}
		b_1 + b_2 = 1\\
		c_2b_2 = \frac 12
	\end{cases}
	.\] 
	La zero stabilità segue dalla Lipschitizanità della dinamica del problema.\\
	\begin{dimo}
		Per $RK2$\\
		\[
			u_{n+1} = u_n = h \sum^{s}_{i=1}b_iK_i \ \ \ \Phi(t_n,u_n) = \sum^{s}_{i=1}b_iK_i = b_1f(t_n,u_n) + b_2f(t_n+c_2h,u_nhc_2f_n)
		.\] 
		Dobbiamo dimostrare che
		\[
			\begin{split}
				&|\Phi(t_n,u_n) - \Phi(t_n,v_n)|\leq \Lambda|u_n-v_n|\\
							       &|\Phi(t_n,u_n) - \Phi(t_n,v_n)|=\\
							       &=|b_1f(t_n,u_n) + b_2f(t_n+c_2h,u_n+hc_2f(t_n,u_n)-b_1f(t_n,v_n)-b_2f(t_n+c_2h,v_n+hc_2f(t_n,v_n)|\\
							       &  \leq |b_1||f(t_n,u_n)-f(t_n,v_n)| + |b_2||f(t_n+c_2h,u_n + hc_2f(t_n,u_n)) - f(t_n+c_2h,v_n+hc_2f(t_n,v_n)|\\
							       & \leq |b_1|L|u_n-v_n|+|b_2|L|u_n+hc_2f(t_n,u_n)-v_n-hc_2f(t_n,u_n)|\\
							       &\leq |b_1|L|u_n-v_n| + |b_2|L(|u_n-v_n|+|hc_2||f(t_n,u_n)-f(t_n,v_n)|)\\
							       &\leq |b_1|L|u_n-v_n|+ |b_2|L|u_n-v_n| + |b_2|Lh|c_2|L|u_n-v_n|\\
							       &\leq [(|b_1| + |b_2|)L + |c_2b_2|L^2h]|u_n-v_n|\\
							       &\leq [(|b_1| + |b_2|)L + |c_2b_2|L^2h_0]|u_n-v_n| \leq \Lambda |u_n-v_n| \ \ \forall h\leq h_0
			\end{split}
		.\] 
	\end{dimoo}
	\subsection{RK4}
	$u_{n+1} = u_n + h(\frac 16 K_1 + \frac 13 K_2 + \frac 13 K_3 + \frac 14 K_4$\\
	Con:\\
	$K_1 = f(t_n,u_n)$\\
	$K_2 = f(t_n + \frac h 2, u_n + \frac h 2 K_1)$\\
	$K_3 = f(t_n + \frac h 2, u_n + \frac h 2 K_2)$\\
	$K_4 = f(t_n + \frac h 2, u_n + \frac h 2 K_3)$\\
	\textbf{Proprietà}
\begin{itemize}
	\item  L'ordine di convergenza di un metodo RK a $s$ stadi non può essere maggiore di $s$
	\item Non esistono metodi espliciti a  $s$ stadi di oridne $s$ se $s\geq 5$
\end{itemize}
\subsection{Assoluta stabilità}

\begin{cases}
$	\dot y = \lambda y\\
	y(0) = 1$
\end{cases}\ \ \ $Re\lambda < 0 \Rightarrow  |y(t)| = |e^{Re(\lambda)e^{iIm(\lambda)}| \xrightarrow{t \rightarrow \infty} 0$\\
	$u_n \xrightarrow{n \rightarrow \infty} 0$  con $h$ fissato
	\[
	\begin{split}
		K_i &= f(t_n + c_ih,u_n + h \sum^{s}_{k=1}a_{ij}K_j) = \lambda (u_n + h \sum^{s}_{j-1}a_{ij}K_j)\\
		u_{n+1} &= u_n + h \sum^{s}_{i=j}b_iK_i
	\end{split}
	.\] 
	\textbf{Notazione }\\
	\[
		\K = (K_1,\ldots,K_s)^T\ \ \ \mathbb{1} = (1,\ldots, 1)^T \ \ A = \{a_{ij}\}_{i,j = 1,\ldots,s}
	.\] 
	 \[
	\begin{split}
		\K = \lambda u_n\mathbb 1 + h\lambda A\K & \Leftrightarrow(I - zA)\K = \lambda u_n\mathbb 1\\
							 & \Leftrightarrow \K = (I-zA)^{-1}\lambda u_n\mathbb 1
	\end{split}
	.\] 
	con $z = h\lambda$\\
	$
		b = (b_1,b_2,\ldots,b_s)^T \Rightarrow  u_{n+1}= u_n + hb^T\K
	$\\
	$u_{n+1} = u_n + hb^T(I-zA)^{-1}\lambda u_n\mathbb 1$\\
	$u_{n+1}= (1 + zb^T(I-zA)^{-1}\mathbb 1)u_n= R(z)u_n$ dove  $R(z)$ è la funzione di stabilità\\
	$u_{n+1} = R(z)u_n = (R(z))^{n+1}u_0 \Rightarrow  u_n \xrightarrow{n \rightarrow +\infty} 0 \Leftrightarrow |R(z)| \leq 1$ \\
	\[
	\begin{cases}
		u_{n+1} - hb^T\K = a_n\\
		(I-zA)\K = \lambda u_n\mathbb 1
		\end{cases} \Leftrightarrow \left( \begin{array}{c|c}
			1 &-hb^T\\
			\hline
			0 & I-zA
\end{array} \right) \matrice{u_{n+1}\\\K } = u_N\matrice{1\\\lambda \mathbb 1}
	\]
	La soluzione si puo calcolare con il metodo di Cramer\\
	\begin{center}
	\includegraphics[scale=0.18]{cramer.png}
	\end{center}\\
	Quindi $u_{n+1} = \det(I-zA + Z\mathbb 1 b^T)u_n = R(z)u_n$   $z$ compone linearmente e il determinante produce un polinomio in $z$ \\
	 \[
	\begin{split}
		u_{n+1} &= u_n + hb^T\K = u_n + zb^T\frac \K \lambda\\
			&= u_n + zb^T(u_n\mathbb 1 + h A\K )\\
			&= u_n + zb^T \left(u_n\mathbb 1 + \frac{zA\K}\lambda \right)\\
			&= u_n + zb^T \left(u_n\mathbb 1 + zAu_n \mathbb 1 + \frac{(zA)^2\K}\lambda \right)\\
			&=u_n + zb^T \left(u_n\mathbb 1 + zA u_n \mathbb 1 + \frac{(zA)^2\K}\lambda \right)
	\end{split}
	.\] 
	Poichè $A$ strettamente diagonale inferiore di dimensione $s$ $A$ è nihil potente di oridne $s$ ( $A^s = 0$)\\
	$u_{n+1} = u_n + zb^T[u_n\mathbb 1 (1 + zA + (zA)^2 + \ldots + (zA)^{s-1}]$\\
	$u_{n+1} = u_n(1 + zb^t \sum^{s-1}_{j=0}(zA)^j\mathbb 1) = u_nR(z) \Rightarrow $  $R(z)$ è un polinomio di grado  $s$ in $z$\\
	Se  $s = 1,2,3,4$ allora  $R(z)$ è lo sviluppo di Taylor di $e^z$ all'ordine  $s$, indipendentemente da $A$ e $b$.
\end{document}
